
\subsection{Exemple PDDL}\label{ex_pddl}

Nous appliquons dans cette section notre parser sur une instanciation de problème PDDL issus de \url{https://github.com/arii/planning_exercises/tree/master}. Le domaine correspond à un domaine Block World classique avec les quatre actions classiques de manipulation d'objets (\textit{"pick-up"}, \textit{"put-down"}, \textit{"stack"},\textit{"unstack"}), la définition complète du domaine se trouve dans l'annexe \ref{domain_bw}. Le problème contient deux objets bloc A et bloc B, initialement le bras ne tient rien, le bloc B est posé sur la table et le bloc A au-dessus de B. L'objectif est de poser A sur la table et que rien ne soit posé sur B. Les conditions de l'objectif sont donc respectées dans deux états : A est posé sur la table et B est posé sur la table, ou bien A est posé sur la table et B est posé sur A.

\subsection*{Définition du  domaine}
%     \begin{lstlisting}[language=Lisp,basicstyle=\ttfamily\footnotesize]
% (define (domain bw)
%   (:requirements :strips)
%   (:constants red green blue)
%   (:predicates
%       (on ?x ?y)      ; object ?x is on ?object ?y
%       (on-table ?x)   ; ?x is directly on the table
%       (clear ?x)      ; ?x has nothing on it
%       (arm-empty)     ; robot isn't holding anything
%       (holding ?x))   ; robot is holding ?x

%   ;; les quatres actions classique de manipulation d'objet
%   (:action pick-up
%      :parameters (?ob1)
%      :precondition (and (clear ?ob1) (on-table ?ob1) (arm-empty))
%      :effect
%      (and (not (on-table ?ob1))
%   	     (not (clear ?ob1))
% 	     (not (arm-empty))
% 	     (holding ?ob1)))

%   (:action put-down
%      :parameters (?ob)
%      :precondition (holding ?ob)
%      :effect
%      (and (not (holding ?ob))
% 	   (clear ?ob)
% 	   (arm-empty)
% 	   (on-table ?ob)))

%   (:action stack
%      :parameters (?sob ?sunderob)
%      :precondition (and (holding ?sob) (clear ?sunderob))
%      :effect
%      (and (not (holding ?sob))
% 	   (not (clear ?sunderob))
% 	   (clear ?sob)
% 	   (arm-empty)
% 	   (on ?sob ?sunderob)))

%   (:action unstack
%      :parameters (?sob ?sunderob)
%      :precondition (and (on ?sob ?sunderob) (clear ?sob) (arm-empty))
%      :effect
%      (and (holding ?sob)
% 	   (clear ?sunderob)
% 	   (not (clear ?sob))
% 	   (not (arm-empty))
% 	   (not (on ?sob ?sunderob))))
% )
%     \end{lstlisting}

\subsection*{Définition du problème}

Lié au domaine défini dans l'annexe \ref{domain_bw}, on introduit une définition de problème. Initialement A et B sont sur la table et on cherche à atteindre un état où B se trouve sur A :

\begin{lstlisting}[language=Lisp,basicstyle=\ttfamily\footnotesize]
(define (problem bw-simple)
  (:domain bw)
  (:objects a b )
  (:init (clear a) (arm-empty) (on a b) (on-table b))
  (:goal (and (on-table a) (clear b))))
    \end{lstlisting}




\begin{figure}[h!]
    \[
        \xymatrix @R=0.1em @C=5em {
            *++[o][F-]{S0} \ar@/_/[r]_{unstack\ a\ b} & 
            *++[o][F-]{S1} \ar@/_/[l]_{stack\ a\ b} \ar@/_/[r]_{put\ down\ a} &
            *++[o][F-]{S2} \ar@/_/[l]_{pick\ up\ a} \ar@/_/[r]_{pick\ up\ b} \ar@(r,u)[]_{\varnothing} &
            *++[o][F-]{S3} \ar@/_/[l]_{put\ down\ b} \ar@/_/[r]_{stack\ b\ a} &
            *++[o][F-]{S4} \ar@/_/[l]_{unstack\ b} \ar@(r,u)[]_{\varnothing}
        }
    \]
    \centering
    \captionsetup{width=0.8\textwidth}
    \caption{LTS obtenu pour le problème simple de block world. Avec S0 =\{'clear a', 'on a b', 'on table b', 'arm-empty'\}, S1 =\{'on table b', 'clear b', 'holding a'\}, S2 =\{'on table b','arm empty'', 'clear b','clear a', 'on table a'\}, S3 =\{'holding b','on table a', 'clear a'\}, S4 =\{'on b a', 'on table a', 'arm empty', 'clear b'\}}
    \label{fig:exemple graphe BW}
\end{figure}